\section{讨论与开放问题}\label{sec:discussion}

这个反例受一个简单局部事件支配，而该事件只有在算子平方后才显现。负四边形
通量从 \(A^2\) 中移除相关的奇位移耦合，正通量则以振幅二激活它们。因此，
全负相位是在谱半径平方为八处的自然消去基线。一对正缺陷只有在八点胞中互为
对径时才能留在该屏障之下；其他每种周期八几何都会产生足够的带符号闭步质量
而越过屏障。

该机制解释了猜想为何能在许多小有限阶上成立。目标相位的无限体积谱边略低于
八，而扭曲猜想阈值向八递增。二者首次在 32 阶交叉。此后，同一周期八字在
每个更大的八的倍数上都有效，因为其 Floquet 谱在整个单位圆上一致受控。

任意周期矩恒等式表明，缺陷图景并非周期八的人为产物。任何位于八屏障处或
之下的相位，其缺陷密度都至多为 \(\frac34\)，且相邻与距离二缺陷簇还必须
满足一个更强的加权不等式。这些约束仅为必要条件。低周期前沿说明了原因：
稀疏的近屏障几何可能需要长达 130 的闭步才首次出现正超额量，而五个有界
周期竞争者用端点 Rayleigh 向量排除更为高效。

计算机辅助结果具有两种不同作用。直至 30 阶的有限搜索确立反例的最小性；
有界周期分类确立目标相位附近的结构前沿。两种计算都没有证明无界优化定理。

当前公开状态检查没有在源论文、作者列出的后续材料或 2026 年 8 月 20 日执行
的窄范围直接检索中发现公开的直接解决。投稿前应刷新该评估，且不应把它理解
为绝对优先权声明。

本文留下若干自然问题。

\textbf{问题 1.} 在任意本原周期的周期相位中，目标相位是否全局最优？

\textbf{问题 2.} 是否存在一个包含非周期赋号且目标相位仍然最优的有意义
无限体积类？

\textbf{问题 3.} 哪些不是八的倍数的偶数阶允许反例？这些阶数上的精确有限
最小值是什么？

\textbf{问题 4.} 能否把层级 \(F_k=M_{k+1}-8M_k\) 转化为任意周期上的有限
结构分类，还是每个尺度都会出现真正新的局部模式？

\textbf{问题 5.} 是否存在一种无需有限轨道分类即可解释对径缺陷几何的通量
相变分原理？

这些问题的答案将推动本文的有界结果与必要条件走向无界通量相定理。本文任何
地方都没有假定这些问题已有肯定答案。
